% Provenance (never printed): NSI premiere "Systemes d'exploitation" (role, commandes de base, droits et
% permissions, logiciel libre). See sources/book2/BRIEFS.md.
\chapter{Operating Systems and the Command Line}\label{ch:g11:os-and-command-line}

Two programs want to print at the same moment; two others want to write the same file; a third would like to read a
file that belongs to someone else. None of them may decide alone. A program that runs underneath all the others, and
that alone may touch the hardware, decides for them: the \omterm{def:g11:os-and-command-line:kernel}{kernel} of the operating system.

This chapter looks at what that \omterm{def:g11:os-and-command-line:kernel}{kernel} does, then works with it the way programmers do every day, by typing \omterm{def:g11:os-and-command-line:shell}{commands}:
moving through the tree of files, copying and removing, setting who may read what, and joining small \omterm{def:g11:os-and-command-line:shell}{commands} into
pipelines that answer questions about the town's data.

\begin{recall}
A program is a list of instructions run by the computer; the operating system is the program that starts the
others and shares the machine between them (Book 1, grade 4). Files are kept in folders, folders in other folders,
and a path names a file by the folders that lead to it (Book 1, grade 3). A text file holds characters only, line by
line (Book 1, grade 8).
\end{recall}

\section{What the kernel does}\label{sec:g11:os-and-command-line:kernel}

\begin{definition}[Kernel, system call]\label{def:g11:os-and-command-line:kernel}
The \emph{kernel}\index{kernel} is the part of the operating system that runs with full control of the hardware. It
shares the processor between the running programs, gives each its own memory, keeps the files, drives the devices,
and checks who may do what. A program asks the kernel for a service through a \emph{system call}\index{system call},
such as opening a file, reading from it, or starting another program.
\end{definition}

A program never drives the disk or the printer itself. When Python runs \code{open("loans.csv")}, it ends in a \omterm{def:g11:os-and-command-line:kernel}{system
call}; the \omterm{def:g11:os-and-command-line:kernel}{kernel} checks that the file exists and that this \omterm{def:g11:os-and-command-line:permission}{user} may read it, and only then hands back a number by
which the program will refer to the open file. Two programs that print at once get their pages one after the other,
because both go through the \omterm{def:g11:os-and-command-line:kernel}{kernel}, which keeps a queue.

\begin{omfigure}
\begin{tikzpicture}[font=\small, l/.style={draw, text width=92mm, minimum height=8mm, align=center}]
  \node[l, fill=black!8] (hw) at (0,0) {hardware: processor, memory, disks, network, screen, printer};
  \node[l, fill=omDef!15, above=0mm of hw] (k) {kernel: processes, memory, files, devices, rights};
  \node[l, fill=omDef!6, above=0mm of k] (s) {system programs: the shell, \code{ls}, \code{cp}, \code{sort}, \dots};
  \node[l, fill=omMeth!8, above=0mm of s] (a) {applications: a browser, an editor, your Python programs};
  \node[above=2mm of a] (u) {the user};
  \draw[<->, thick, omThm] ([xshift=48mm]a.north) -- node[right, align=left, font=\footnotesize] {system\\calls}
    ([xshift=48mm]k.north);
\end{tikzpicture}

{\small The layers of a computer. Programs above the \omterm{def:g11:os-and-command-line:kernel}{kernel} reach the hardware only by asking the \omterm{def:g11:os-and-command-line:kernel}{kernel}, through
\omterm{def:g11:os-and-command-line:kernel}{system calls}.}
\end{omfigure}

The operating system is the \omterm{def:g11:os-and-command-line:kernel}{kernel} and the system programs around it. Some are \omterm{def:g11:os-and-command-line:free-software}{free software}, some are not.

\begin{definition}[Free software]\label{def:g11:os-and-command-line:free-software}
A program is \emph{free software}\index{free software} when its licence gives every \omterm{def:g11:os-and-command-line:permission}{user} the freedom to share it,
to study how it works, and to modify it, which requires access to its source code. A \emph{proprietary} program's
licence withholds some of these freedoms.
\end{definition}
% ledger: g11:free-software

A related term, \textit{open source}, comes with its own definition, a list of ten criteria starting with free
redistribution and access to the source code; in practice most licences meet both definitions.
% ledger: g11:osd-10

\begin{dated}{2026}
The operating systems in common use form a few families. The Linux \omterm{def:g11:os-and-command-line:kernel}{kernel}, \omterm{def:g11:os-and-command-line:free-software}{free software}, runs most servers and,
modified, Android phones, whose \omterm{def:g11:os-and-command-line:kernel}{kernel} is based on a long-term-supported Linux \omterm{def:g11:os-and-command-line:kernel}{kernel}. Apple's systems run its
computers and phones. Microsoft's Windows runs many personal computers. They differ in many details, but the
\omterm{def:g11:os-and-command-line:shell}{commands} of this chapter work on Linux and on Apple's computers, and on Windows inside a Linux environment.
% ledger: g11:android-linux
\end{dated}

\section{The file-system tree}\label{sec:g11:os-and-command-line:tree}

\begin{definition}[File-system tree, root directory]\label{def:g11:os-and-command-line:fs-tree}
The \emph{file-system tree}\index{file-system tree} is the arrangement of all the directories (folders) and files of
a system: each directory may contain files and other directories, and all of them are inside one directory, the
\emph{root directory}\index{root directory}, written \code{/}.
\end{definition}

\begin{definition}[Absolute path, relative path, working directory]\label{def:g11:os-and-command-line:abs-rel}
An \emph{absolute path}\index{absolute path} names a file from the root, starting with \code{/}, such as
\code{/home/ada/library/books/books.csv}. Each running program has a \emph{working directory}\index{working
directory}, and a \emph{relative path}\index{relative path}, which does not start with \code{/}, names a file from
there. In a path, \code{.} is the directory itself and \code{..} its parent.
\end{definition}

The \omterm{def:g11:os-and-command-line:shell}{shell} is the program that reads \omterm{def:g11:os-and-command-line:shell}{commands} and runs them. The sessions of this chapter are captured from a real
\omterm{def:g11:os-and-command-line:shell}{shell}: each line after \texttt{\$} is a \omterm{def:g11:os-and-command-line:shell}{command}, and the lines below it are what it printed. They run in a fresh
directory, shown as \texttt{\textasciitilde/session}, where \texttt{\textasciitilde} is the \omterm{def:g11:os-and-command-line:permission}{user}'s home directory.

\omoutput{code/grade-11/03-os-and-command-line/outputs/tree_moves.txt}{Making directories and moving through them.
When its output goes to a file, as here, \code{ls} prints one name per line; at a terminal it uses
columns.}\label{out:g11:os-and-command-line:tree-moves}

\code{mkdir} makes a directory, \code{cd} changes the \omterm{def:g11:os-and-command-line:abs-rel}{working directory}, \code{pwd} prints it. A name starting with a
dot is hidden: \code{ls} leaves it out, \code{ls -a} shows it, along with \code{.} and \code{..}. Settings files are
hidden this way, so that they do not clutter the listing.

\omoutput{code/grade-11/03-os-and-command-line/outputs/paths.txt}{\omterm{def:g11:os-and-command-line:abs-rel}{Relative paths} with \code{..}, and the \omterm{def:g11:os-and-command-line:abs-rel}{absolute path}
that \code{realpath} computes.}\label{out:g11:os-and-command-line:paths}

\begin{omfigure}
\begin{tikzpicture}[font=\small\ttfamily, level distance=11mm, sibling distance=26mm,
    every node/.style={draw, rounded corners=1.5pt, inner sep=2pt, fill=white},
    edge from parent/.style={draw, black!60}]
  \node {\textasciitilde/session}
    child {node (lib) {library}
      child {node (bk) {books}
        child {node[fill=omMeth!12] (csv) {books.csv}}}
      child {node[fill=omDef!15, very thick] (mb) {members}}};
  \begin{scope}[every node/.style={font=\footnotesize\rmfamily, text=omThm}]
    \draw[->, thick, omThm] (mb.east) to[bend right=50] node[right] {1: \code{..}} (lib.east);
    \draw[->, thick, omThm] (lib.west) to[bend right=50] node[left] {2: \code{books}} (bk.west);
    \draw[->, thick, omThm] (bk.west) to[bend right=50] node[left] {3: \code{books.csv}} (csv.west);
  \end{scope}
  \node[draw=none, fill=none, below=1mm of mb, font=\footnotesize\rmfamily] {working directory};
\end{tikzpicture}

{\small The session's tree. From the \omterm{def:g11:os-and-command-line:abs-rel}{working directory} \code{members}, the \omterm{def:g11:os-and-command-line:abs-rel}{relative path} \code{../books/books.csv}
takes three steps: up to \code{library}, down into \code{books}, then to the file.}
\end{omfigure}

\section{The shell}\label{sec:g11:os-and-command-line:shell}

\begin{definition}[Shell, command, option]\label{def:g11:os-and-command-line:shell}
A \emph{shell}\index{shell} is a program that reads command lines and runs them. A \emph{command}\index{command}
line is a program's name followed by its arguments, separated by spaces. An argument that starts with \code{-} is an
\emph{option}\index{option}: it changes how the command works, such as \code{-a} for \code{ls}.
\end{definition}

\omoutput{code/grade-11/03-os-and-command-line/outputs/copy_move.txt}{\code{cp} copies, \code{mv} moves or renames,
\code{rm} removes a file, \code{rmdir} an empty directory.}\label{out:g11:os-and-command-line:copy-move}

\begin{safety}
\code{rm} does not move a file to a bin: the file is gone, and there is no \omterm{def:g11:os-and-command-line:shell}{command} to bring it back. \code{rm -r}
removes a directory with everything inside it. Before removing, list what will go with \code{ls}, never remove with
a pattern you have not checked, and keep copies of what matters somewhere else.
\end{safety}

Every \omterm{def:g11:os-and-command-line:shell}{command} describes itself. \code{--help} prints a short summary, and \code{man} followed by the \omterm{def:g11:os-and-command-line:shell}{command}'s name
shows its full manual.

\omoutput{code/grade-11/03-os-and-command-line/outputs/manual.txt}{The beginning of \code{ls}'s own summary, cut
by \code{head}, which keeps the first lines.}\label{out:g11:os-and-command-line:manual}

\section{Users, rights and permissions}\label{sec:g11:os-and-command-line:permissions}

\begin{definition}[User, group, permission]\label{def:g11:os-and-command-line:permission}
Each person or service on a system is a \emph{user}\index{user}, with a name and a number, and belongs to one or
more \emph{groups}\index{group}. Each file has an owner and a group, and nine \emph{permission}\index{permission} bits:
read, write and execute, for the owner, for the members of the group, and for all the others. For a directory, read
lets one list it, write lets one create or remove names in it, and execute lets one enter it and reach what is
inside.
\end{definition}

The \omterm{def:g11:os-and-command-line:kernel}{kernel} checks the bits of one class only: the owner's if the \omterm{def:g11:os-and-command-line:permission}{user} owns the file, otherwise the \omterm{def:g11:os-and-command-line:permission}{group}'s if the
\omterm{def:g11:os-and-command-line:permission}{user} is in the file's \omterm{def:g11:os-and-command-line:permission}{group}, otherwise the others'. \code{ls -l} prints them as ten characters: the type (\code{d}
for a directory, \code{-} for a file), then \code{r}, \code{w}, \code{x} or \code{-} for each class.
% ledger: g11:ls-mode

\omoutput{code/grade-11/03-os-and-command-line/outputs/perms.txt}{\code{chmod} changes the bits. Once its owner has
no \code{r}, even the owner cannot read the file, until the bit comes back.}\label{out:g11:os-and-command-line:perms}

\code{chmod} takes the bits either as letters, \code{u}, \code{g}, \code{o} for the classes, \code{+} or \code{-}
and the rights, or as a number in octal: one digit per class, the sum of $4$ for read, $2$ for write and $1$ for
execute.
% ledger: g11:chmod-octal

\begin{omfigure}
\begin{bytefield}[bitwidth=3.4em]{10}
  \bitheader[endianness=big, lsb=0]{0-9} \\
  \bitbox{1}{\code{-}} & \bitbox{1}{\code{r}} & \bitbox{1}{\code{w}} & \bitbox{1}{\code{-}} &
  \bitbox{1}{\code{r}} & \bitbox{1}{\code{-}} & \bitbox{1}{\code{-}} & \bitbox{1}{\code{-}} & \bitbox{1}{\code{-}} &
  \bitbox{1}{\code{-}} \\
  \bitbox{1}{type} & \bitbox{3}{owner: $4 + 2 = 6$} & \bitbox{3}{group: $4$} & \bitbox{3}{others: $0$}
\end{bytefield}

{\small The mode \code{-rw-r-----} of \code{books.csv} after \code{chmod 640}: the owner reads and writes, the \omterm{def:g11:os-and-command-line:permission}{group}
reads, the others have nothing.}
\end{omfigure}

\begin{inpractice}[The administrator]
Every system has an administrator account, called \code{root} on Linux, for which the \omterm{def:g11:os-and-command-line:kernel}{kernel} skips the \omterm{def:g11:os-and-command-line:permission}{permission}
checks. A program run as the administrator can read, change or remove any file; a bug in it, or an attacker who takes
it over, can too. Programs run as an ordinary \omterm{def:g11:os-and-command-line:permission}{user}, with the fewest rights that let them work, and the administrator
account is used only for the tasks that need it.
\end{inpractice}

\section{Redirection and pipes}\label{sec:g11:os-and-command-line:pipes}

\begin{definition}[Standard input and output, redirection, pipe]\label{def:g11:os-and-command-line:redirection}
Each program starts with a \emph{standard input}\index{standard input/output}, where it reads, and a \emph{standard
output}, where it writes; at a terminal they are the keyboard and the screen. A \emph{redirection}\index{redirection}
connects them to a file instead: \code{> f} writes the output into \code{f}, emptying it first, \code{>> f} adds at
the end of \code{f}, and \code{< f} reads the input from \code{f}. A \emph{pipe}\index{pipe}, \code{a | b},
connects the standard output of \code{a} to the standard input of \code{b}.
\end{definition}

\omoutput{code/grade-11/03-os-and-command-line/outputs/redirect.txt}{\code{>} replaces, \code{>>} adds. The
town's \code{books.csv} has $81$ lines, a header and $80$ books.}\label{out:g11:os-and-command-line:redirect}

Small \omterm{def:g11:os-and-command-line:shell}{commands}, each doing one thing, join into pipelines. \code{cut -d, -f3} keeps the third field of each line,
fields being separated by commas; \code{tail -n +2} drops the first line; \code{sort} sorts the lines;
\code{uniq -c} counts each run of equal lines, which must be adjacent, hence the \code{sort} before it; \code{sort -n} sorts by
numerical value; \code{grep -c} counts the lines that match a pattern.
% ledger: g11:uniq-c
% ledger: g11:sort-n

\omoutput{code/grade-11/03-os-and-command-line/outputs/pipes.txt}{Which members borrowed most? Member $55$, with
$12$ loans. The last \omterm{def:g11:os-and-command-line:shell}{command} counts the loans whose return date is empty: $22$ books are still
out.}\label{out:g11:os-and-command-line:pipes}

\begin{omfigure}
\begin{tikzpicture}[font=\footnotesize\ttfamily, b/.style={draw, rounded corners=2pt, fill=omDef!8, minimum height=8mm,
    inner sep=3pt, text height=1.6ex, text depth=0.4ex}]
  \node[b, fill=omMeth!12] (f) at (0,0) {loans.csv};
  \node[b, right=5mm of f] (c) {cut -d, -f3};
  \node[b, right=5mm of c] (t) {tail -n +2};
  \node[b, right=5mm of t] (s) {sort};
  \node[b, right=5mm of s] (u) {uniq -c};
  \node[b, right=5mm of u] (n) {sort -n};
  \node[b, right=5mm of n] (l) {tail -3};
  \draw[->, thick] (f) -- (c);
  \foreach \a/\b in {c/t, t/s, s/u, u/n, n/l} \draw[->, thick, omThm] (\a) -- (\b);
  \node[below=1mm of c, font=\scriptsize\rmfamily] {member numbers};
  \node[below=1mm of u, font=\scriptsize\rmfamily] {count per member};
  \node[below=1mm of l, font=\scriptsize\rmfamily] {the three largest};
\end{tikzpicture}

{\small A pipeline: each red arrow is a \omterm{def:g11:os-and-command-line:redirection}{pipe} from one program's standard output to the next one's \omterm{def:g11:os-and-command-line:redirection}{standard input}.
The six programs run at the same time, each working on the lines as they arrive.}
\end{omfigure}

\section{Exercises}\label{sec:g11:os-and-command-line:exercises}

\begin{exercise}[$\star$]\label{exo:g11:os-and-command-line:1}
The \omterm{def:g11:os-and-command-line:abs-rel}{working directory} is \code{/home/ada/library/books}. What does \code{pwd} print after \code{cd ../members}?
After \code{cd ../../notes}?
\end{exercise}

\begin{exercise}[$\star$]\label{exo:g11:os-and-command-line:2}
The \omterm{def:g11:os-and-command-line:abs-rel}{working directory} is \code{/home/ada/library/members}. Give the \omterm{def:g11:os-and-command-line:abs-rel}{absolute path} of \code{../books/./books.csv}.
\end{exercise}

\begin{exercise}[$\star$]\label{exo:g11:os-and-command-line:3}
Write the mode \code{rw-r-----} in octal, and \code{chmod 754} in letters.
\end{exercise}

\begin{exercise}[$\star$]\label{exo:g11:os-and-command-line:4}
A directory holds \code{a.txt} and \code{b.txt}. What does \code{ls} print after \code{mv a.txt c.txt}? After
\code{mv b.txt c.txt}?
\end{exercise}

\begin{exercise}[$\star$]\label{exo:g11:os-and-command-line:5}
Write a \omterm{def:g11:os-and-command-line:shell}{command} line that prints the number of books in \code{books.csv}, without counting its header.
\end{exercise}

\begin{exercise}[$\star\star$]\label{exo:g11:os-and-command-line:6}
To add a line at the top of \code{f.txt}, a student types \code{cat top.txt f.txt > f.txt}. What does
\code{f.txt} contain afterwards, and why? Write a \omterm{def:g11:os-and-command-line:shell}{command} line that works.
\end{exercise}

\begin{exercise}[$\star\star$]\label{exo:g11:os-and-command-line:7}
A file has owner \code{ana}, \omterm{def:g11:os-and-command-line:permission}{group} \code{teachers} and mode \code{640}. The \omterm{def:g11:os-and-command-line:permission}{group} has the members \code{ana} and
\code{ben}. Which of \code{ana}, \code{ben} and \code{chloe} may read it? Write it?
\end{exercise}

\begin{exercise}[$\star\star$]\label{exo:g11:os-and-command-line:8}
Write a pipeline that prints the three authors with the most books in \code{books.csv}, with their counts.
\end{exercise}

\begin{exercise}[$\star\star$]\label{exo:g11:os-and-command-line:9}
A directory is named \code{My Notes}. What happens with \code{cd My Notes}? Give two ways to enter it.
\end{exercise}

\begin{exercise}[$\star\star$]\label{exo:g11:os-and-command-line:10}
What does \code{g.txt} contain after \code{echo a >> g.txt} and then \code{echo b >> g.txt}, if it did not exist
before? And if the second \code{>>} is a \code{>}?
\end{exercise}

\begin{exercise}[$\star\star$]\label{exo:g11:os-and-command-line:11}
Give a \code{chmod} \omterm{def:g11:os-and-command-line:shell}{command} that lets the \omterm{def:g11:os-and-command-line:permission}{group} read a file but not write it, without changing the bits of the owner
and of the others.
\end{exercise}

\begin{exercise}[$\star\star\star$]\label{exo:g11:os-and-command-line:12}
How many different values can the nine \omterm{def:g11:os-and-command-line:permission}{permission} bits take? In how many of them may the others write?
\end{exercise}

\begin{exercise}[$\star\star\star$]\label{exo:g11:os-and-command-line:13}
Explain why any \omterm{def:g11:os-and-command-line:abs-rel}{relative path}, together with the \omterm{def:g11:os-and-command-line:abs-rel}{absolute path} of the \omterm{def:g11:os-and-command-line:abs-rel}{working directory}, can be written as an
\omterm{def:g11:os-and-command-line:abs-rel}{absolute path} that contains neither \code{.} nor \code{..}. Describe the steps, and why they always end.
\end{exercise}

\begin{exercise}[$\star\star\star$]\label{exo:g11:os-and-command-line:14}
Write a \omterm{def:g11:os-and-command-line:shell}{shell} script \code{count_loans.sh} that takes a member number as its argument, written \code{$1} inside the
script, and prints how many loans of \code{loans.csv} that member made. It should print a usage message and stop
with an error status if it gets no argument.
\end{exercise}

\begin{exercise}[$\star\star\star$]\label{exo:g11:os-and-command-line:15}
To find the largest book number in \code{loans.csv}, a student types
\smallskip\noindent\hspace*{2em}\code{cut -d, -f2 loans.csv | tail -n +2 | sort | tail -1}

\smallskip\noindent and gets \code{9}. Why? Fix the \omterm{def:g11:os-and-command-line:shell}{command} line.
\end{exercise}

\section{Problem: Who May Read the Timetable?}\label{sec:g11:os-and-command-line:problem}

\begin{problem}[{Weekend problem --- the fewest rights that meet a policy}]\label{pb:g11:os-and-command-line:1}
A school's small server has four \omterm{def:g11:os-and-command-line:permission}{users}: the teachers \code{ana} and \code{ben}, in the \omterm{def:g11:os-and-command-line:permission}{group} \code{teachers}, and the
students \code{chloe} and \code{dev}. The program below holds its files with their owner, \omterm{def:g11:os-and-command-line:permission}{group} and mode, and computes
who may read and write each one, following the \omterm{def:g11:os-and-command-line:kernel}{kernel}'s rule, including the execute bit of every directory above.

\omcode{code/grade-11/03-os-and-command-line/listings/timetable.py}{1}{38}{The server's files and the \omterm{def:g11:os-and-command-line:kernel}{kernel}'s rule,
simulated.}\label{lst:g11:os-and-command-line:timetable}

\omoutput{code/grade-11/03-os-and-command-line/outputs/timetable.txt}{Who may read (\code{r}) and write (\code{w})
each file, with its present mode.}\label{out:g11:os-and-command-line:timetable}

The head of the school then writes a policy: everyone reads the timetable and the news, but only teachers change the
timetable and only \code{ben} the news; only teachers may see the marks; the exam is \code{ben}'s alone. The second
program tries all $512$ modes of each file, parents first, and keeps the one with the fewest bits among those that
meet the policy exactly.

\omcode{code/grade-11/03-os-and-command-line/listings/policy.py}{1}{35}{The policy, and a search for the smallest
modes that meet it.}\label{lst:g11:os-and-command-line:policy}

\omoutput{code/grade-11/03-os-and-command-line/outputs/policy.txt}{The modes to set, and the bits that
change.}\label{out:g11:os-and-command-line:policy}

\medskip\noindent\textbf{Part I --- Users and groups.}
\begin{enumerate}
  \item For the file \code{/school/news.txt}, which class of bits applies to each of the four \omterm{def:g11:os-and-command-line:permission}{users}?
  \item Why does \code{bits} test the owner before the \omterm{def:g11:os-and-command-line:permission}{group}?
  \item \code{ana} owns \code{/school/news.txt}'s directory but not the file. Can she change the file? Can she remove
    it?
  \item What does \code{can} check about the directories above the file, and why?
\end{enumerate}

\medskip\noindent\textbf{Part II --- The present modes.}
\begin{enumerate}[resume]
  \item Write the mode of each file as ten characters, as \code{ls -l} would.
  \item Explain the row of \code{/school/timetable.txt}: why may \code{chloe} change it?
  \item Explain why \code{dev} may read the marks.
  \item \code{ana} owns the timetable; how could a student who may write it harm her?
\end{enumerate}

\medskip\noindent\textbf{Part III --- Directories.}
\begin{enumerate}[resume]
  \item What do the bits $r$, $w$, $x$ allow on \code{/school/marks}?
  \item If \code{/school/marks} had mode \code{774}, could \code{chloe} list its names? Read \code{g11.csv}?
  \item If it had mode \code{771}, could she read \code{g11.csv} if she knew its name?
  \item With the new mode \code{770}, why do $8$ modes of \code{g11.csv} meet the policy?
  \item Why does \code{policy.py} handle the directories before the files inside them?
\end{enumerate}

\medskip\noindent\textbf{Part IV --- The least privileges.}
\begin{enumerate}[resume]
  \item Check by hand that \code{664} meets the policy for the timetable, and that no mode with fewer bits does.
  \item For each file, write the \code{chmod} \omterm{def:g11:os-and-command-line:shell}{command}, in octal and in letters.
  \item Count the bits changed for each file from its old and new modes, and check the total.
  \item The exam file's \omterm{def:g11:os-and-command-line:permission}{group} is \code{teachers}. Why is \code{600} still needed although \code{ana} is a teacher?
  \item A new teacher arrives. Which single change gives her every right the policy grants to teachers, without any \code{chmod}?
  \item Why is it better to give the fewest rights that work, rather than open everything and close what goes wrong?
  \item State the result: the modes to set and the number of \omterm{def:g11:os-and-command-line:permission}{permission} bits that change.
\end{enumerate}
\end{problem}
